% Options for packages loaded elsewhere
\PassOptionsToPackage{unicode}{hyperref}
\PassOptionsToPackage{hyphens}{url}
%
\documentclass[
]{article}
\usepackage{amsmath,amssymb}
\usepackage{iftex}
\ifPDFTeX
  \usepackage[T1]{fontenc}
  \usepackage[utf8]{inputenc}
  \usepackage{textcomp}
\else
  \usepackage{unicode-math}
  \defaultfontfeatures{Scale=MatchLowercase}
  \defaultfontfeatures[\rmfamily]{Ligatures=TeX,Scale=1}
\fi
\usepackage{lmodern}
\IfFileExists{upquote.sty}{\usepackage{upquote}}{}
\IfFileExists{microtype.sty}{\usepackage[]{microtype}\UseMicrotypeSet[protrusion]{basicmath}}{}
\makeatletter
\@ifundefined{KOMAClassName}{\IfFileExists{parskip.sty}{\usepackage{parskip}}{\setlength{\parindent}{0pt}\setlength{\parskip}{6pt plus 2pt minus 1pt}}}{\KOMAoptions{parskip=half}}
\makeatother
\usepackage{xcolor}
\setlength{\emergencystretch}{3em}
\providecommand{\tightlist}{\setlength{\itemsep}{0pt}\setlength{\parskip}{0pt}}
\setcounter{secnumdepth}{-\maxdimen}
\ifLuaTeX\usepackage{selnolig}\fi
\usepackage{bookmark}
\IfFileExists{xurl.sty}{\usepackage{xurl}}{}
\urlstyle{same}
\hypersetup{hidelinks,pdfcreator={LaTeX via pandoc}}
\author{}
\date{}
\begin{document}

\section{Separating Patterns in Finite Groups: Exact Families and Small Computations}
\textbf{Igor Tomilson}\\
Independent researcher

\subsection{Abstract}
Let \(G\) be a finite group, let \(q\ge 2\), and let \(f:G\to[q]\). For a finite set \(Y\subseteq G\), consider the translated observation map
\[
\Phi_{f,Y}(x)=(f(xy))_{y\in Y}.
\]
The separating number \(\operatorname{sep}_q(G)\) is the minimum size of \(Y\) for which \(\Phi_{f,Y}\) is injective for some \(f\).

This note records exact infinite families and finite computational results. First, for every cyclic group \(C_N\),
\[
\operatorname{sep}_q(C_N)=\lceil \log_q N\rceil.
\]
Second, for the elementary abelian binary groups \(V_d=(C_2)^d\),
\[
\operatorname{sep}_2(V_d)=
\begin{cases}
3,&d=2,\\
5,&d=4,\\
d,&\text{otherwise}.
\end{cases}
\]
More generally, for \(E_n(p)=(C_p)^n\) and an alphabet of size \(p^r\), we give a complete classification. Finally, we give the exact computational result
\[
\operatorname{sep}_2((C_3)^3)=6>\lceil\log_2 27\rceil=5,
\]
as well as the sharp values
\[
\operatorname{sep}_2((C_3)^4)=7,
\qquad
\operatorname{sep}_2((C_5)^3)=7,
\]
and explicit certificates showing that every one of the 51 groups of order \(32\) satisfies
\[
\operatorname{sep}_2(G)=5.
\]
All computational results are accompanied by reproducible scripts and certificates.

\subsection{1. Introduction}
Kang and Hsieh introduced a separating-pattern parameter for finite groups in their study of information and locality in Cayley graphs.

Let \(G\) be a finite group and let \(f:G\to[q]\). For an ordered finite set
\[
Y=(y_1,\ldots,y_t)\subseteq G
\]
define
\[
\Phi_{f,Y}(x)=\bigl(f(xy_1),\ldots,f(xy_t)\bigr).
\]
We call \((f,Y)\) separating if \(\Phi_{f,Y}\) is injective, and define
\[
\operatorname{sep}_q(G)=\min\{|Y|:\exists f:G\to[q]\text{ such that }(f,Y)\text{ is separating}\}.
\]
A counting argument gives
\[
\operatorname{sep}_q(G)\ge\left\lceil\log_q|G|\right\rceil.
\]
The reproducibility package is available at \url{https://github.com/spookyyST/separating-numbers-finite-groups}.

\subsection{2. Cyclic groups}
\subsubsection{Theorem 2.1}
For every \(q\ge2\) and \(N\ge2\),
\[
\boxed{\operatorname{sep}_q(C_N)=\lceil\log_qN\rceil.}
\]
\subsubsection{Proof}
Let \(r=\lceil\log_qN\rceil\). The counting bound gives \(\operatorname{sep}_q(C_N)\ge r\).

Write \(C_N=\langle s\rangle\). Since \(N\le q^r\), choose a cyclic \(q\)-ary word \(c=(c_0,c_1,\ldots,c_{N-1})\) of length \(N\) whose cyclic factors of length \(r\) are pairwise distinct. Such arbitrary-length cut-down de Bruijn sequences exist whenever \(N\le q^r\).

Define \(f(s^i)=c_i\) and take \(Y=(1,s,s^2,\ldots,s^{r-1})\). Then
\[
\Phi_{f,Y}(s^i)=(c_i,c_{i+1},\ldots,c_{i+r-1}),
\]
with indices modulo \(N\). These \(N\) words are pairwise distinct, so \(\operatorname{sep}_q(C_N)\le r\). Together with the lower bound, \(\operatorname{sep}_q(C_N)=r\). \(\square\)

\subsubsection{Remark 2.2}
Every cyclic group attains the information lower bound for every alphabet size \(q\ge2\).

\subsection{3. Elementary abelian binary groups}
Let \(V_d=\mathbb F_2^d\).

\subsubsection{Theorem 3.1}
For every \(d\ge1\),
\[
\boxed{\operatorname{sep}_2(V_d)=
\begin{cases}
3,&d=2,\\
5,&d=4,\\
d,&\text{otherwise}.
\end{cases}}
\]
Since \(|V_d|=2^d\), \(\operatorname{sep}_2(V_d)\ge d\).

\subsubsection{3.1 Odd dimensions}
Let \(d=2m+1\). Define
\[
f(x)=x_1+\sum_{j=1}^{m}x_{2j}x_{2j+1}
\]
and take \(Y=(0,e_2,e_3,\ldots,e_{2m+1})\). For every \(j\),
\[
f(x+e_{2j})+f(x)=x_{2j+1},\qquad f(x+e_{2j+1})+f(x)=x_{2j}.
\]
Thus the observation word recovers \(x_2,\ldots,x_{2m+1}\), and then
\[
x_1=f(x)+\sum_{j=1}^{m}x_{2j}x_{2j+1}.
\]
Hence \(\operatorname{sep}_2(V_d)=d\) for every odd \(d\).

\subsubsection{3.2 Two-dimensional lifting lemma}
\textbf{Lemma 3.2.} Suppose \((f,Y)\) separates \(V_d\), with \(|Y|=d\) and \(0\in Y\). Then \(\operatorname{sep}_2(V_{d+2})=d+2\).

\textit{Proof.} Write \(V_{d+2}=V_d\times\mathbb F_2^2\) with new coordinates \(a,b\). Define \(F(x,a,b)=f(x)+ab\) and take
\[
Y'=\{(y,0,0):y\in Y\}\cup\{(0,1,0),(0,0,1)\}.
\]
The two new observations determine
\[
F(x,a+1,b)+F(x,a,b)=b,
\]
and
\[
F(x,a,b+1)+F(x,a,b)=a.
\]
Thus \(a,b\) are recovered. After subtracting the known term \(ab\) from the embedded old observations, one recovers the original \(Y\)-word, hence \(x\). Therefore \(Y'\) separates \(V_{d+2}\). The counting bound gives equality. \(\square\)

\subsubsection{3.3 Dimension six}
For \(V_6\), take \(Y=(0,e_1,e_2,e_3,e_4,e_5)\) and
\[
\begin{aligned}
f={}&x_2+x_1x_2+x_1x_3+x_2x_3+x_3x_4+x_2x_5\\
&+x_1x_4x_5+x_3x_4x_5+x_6+x_2x_6+x_3x_6+x_4x_6\\
&+x_1x_4x_6+x_3x_4x_6+x_4x_5x_6.
\end{aligned}
\]
The accompanying verifier checks that
\[
x\longmapsto\bigl(f(x),f(x+e_1),\ldots,f(x+e_5)\bigr)
\]
is a permutation of \(\mathbb F_2^6\). Hence \(\operatorname{sep}_2(V_6)=6\). The lifting lemma gives \(\operatorname{sep}_2(V_{2m})=2m\) for every \(m\ge3\).

\subsubsection{3.4 Exceptional dimensions}
Finite exhaustive verification gives
\[
\operatorname{sep}_2(V_2)=3,\qquad \operatorname{sep}_2(V_4)=5.
\]
For \(d=2\), every normalized two-point pattern is equivalent to \((0,e_1)\). For \(d=4\), every normalized four-point pattern is affine-equivalent to one of two types,
\[
(0,e_1,e_2,e_1+e_2)
\]
or
\[
(0,e_1,e_2,e_3).
\]
The verifier exhaustively checks the corresponding Boolean functions and finds no separating window of size \(d\). Matching upper bounds are also verified. This completes the proof of Theorem 3.1.

\subsection{4. Prime-power alphabets on elementary abelian groups}
Let
\[
E_n(p)=(C_p)^n\cong\mathbb F_p^n,\qquad q=p^r.
\]
Identify the alphabet of size \(q\) with \(\mathbb F_p^r\).

\subsubsection{Theorem 4.1}
For primes \(p\) and integers \(n,r\ge1\),
\[
\operatorname{sep}_{p^r}(E_n(p))=
\begin{cases}
\lceil n/r\rceil,&p\text{ odd},\\
3,&p=2,\ r\ge2,\ n=2r,\\
\lceil n/r\rceil,&p=2,\ r\ge2,\ n\ne2r.
\end{cases}
\]
For \(p=2,r=1\), Theorem 3.1 gives the two exceptional dimensions \(n=2,4\).

\subsubsection{Proof}
The counting bound is
\[
\operatorname{sep}_{p^r}(E_n(p))\ge\lceil n/r\rceil.
\]
Write \(t=\lceil n/r\rceil\). Suppose first that \(p\) is odd. The case
\(t=1\) follows by choosing an injective colouring. Otherwise write a state as
\((u;z)\), where \(u\in\mathbb F_p^r\), and partition the remaining \(n-r\)
coordinates into \(t-1\) nonempty blocks of size at most \(r\). Put the
coordinates of block \(i\) into slots \(z_{i,j}\), treating absent slots as
zero, and define
\[
f_j(u;z)=u_j+\sum_i z_{i,j}^2\qquad(1\le j\le r).
\]
For each block use the shift adding \(1\) to every coordinate in that block.
Its translated difference in component \(j\) is \(2z_{i,j}+1\), so all
\(z\)-coordinates are recovered because \(2\) is invertible. The base colour
then recovers \(u\). Thus \(t\) observations suffice.

Now let \(p=2\) and \(r\ge2\). If \(t=1\), use an injective colouring. If
\(t=2\) and \(r<n<2r\), a two-point window partitions \(E_n(2)\) into
\(2^{n-1}\) pairs. There are \(\binom{2^r}{2}\) unordered pairs of distinct
colours, enough to assign a different colour pair to every orbit. If \(n=2r\),
there are \(2^{2r-1}\) orbits, strictly more than
\(\binom{2^r}{2}\), so two observations are impossible. Identifying
\(E_{2r}(2)\) with the additive group of \(\mathbb F_{2^r}^2\), the colouring
\(f(a,b)=ab\) and window \(\{0,e_1,e_2\}\) recover \(b\) and \(a\) from the
two differences. Hence the exact value is three.

It remains to consider \(t\ge3\). Put \(s=t-1\), \(m=n-r\), and write a
state as \((u,z)\in\mathbb F_2^r\times\mathbb F_2^m\). Choose coordinate
directions \(a_1,\ldots,a_s\) in \(z\) as the nonzero window shifts. The
\(rs\) derivative slots \(D_{a_i}Q_j\) of an \(r\)-bit quadratic map \(Q\)
will recover \(z\). If \(s\) is even, take
\[
Q_1=a_1a_2+a_3a_4+\cdots+a_{s-1}a_s.
\]
If \(s\) is odd, use this pairing through \(a_{s-1}\) in \(Q_1\), and add
\(Q_2=a_1a_s\). These reserve exactly \(s\) derivative slots and recover all
\(a_i\). There are \(rs-s\) remaining slots, while the remaining \(m-s\)
coordinates of \(z\) satisfy \(m-s\le rs-s\). Assign each such coordinate
\(w\) to a distinct unreserved slot \((i,j)\) and add \(a_iw\) to \(Q_j\).
Its derivative gives \(w\), possibly plus an already recovered \(a_i\). Thus
all of \(z\) is recovered. Finally
\[
f(u,z)=u+Q(z)
\]
with window \(\{0,a_1,\ldots,a_s\}\) separates the group: the differences
recover \(z\), and the base colour recovers \(u\). This attains the counting
bound and completes the proof. \(\square\)

\subsection{5. Binary computations in odd characteristic}
\subsubsection{Theorem 5.1}
For \(G=(C_3)^3\cong\mathbb F_3^3\),
\[
\boxed{\operatorname{sep}_2(G)=6.}
\]
Since \(\lceil\log_2|G|\rceil=\lceil\log_2 27\rceil=5\), this group does not attain the information lower bound.

\subsubsection{Computational proof}
Any five-element window can be translated so that it contains the identity. It is therefore enough to consider \(\binom{26}{4}=14,950\) normalized five-windows.

The group \(GL(3,3)\) acts on these windows. The supplied program generates this action using coordinate swaps, nonzero coordinate scalings, and elementary transvections. The 14,950 windows split into ten orbits.

For a fixed representative \(Y\), introduce one Boolean variable \(X_z\) for each \(z\in G\). For every pair \(g\ne h\), separation requires
\[
\bigvee_{y\in Y}\bigl(X_{g+y}\operatorname{xor}X_{h+y}\bigr).
\]
The program converts these conditions to CNF with auxiliary XOR variables and solves every orbit representative exactly. All ten instances are UNSAT.

Thus no separating binary window of size \(5\) exists. The file \texttt{c3xc3xc3\_size6\_certificate.json} contains an explicit six-element window and Boolean colouring. A direct verifier checks that its 27 translated words are pairwise distinct. Therefore \(\operatorname{sep}_2((C_3)^3)=6\). \(\square\)

\subsubsection{Theorem 5.2}
The following two groups attain their binary counting bounds:
\[
\boxed{\operatorname{sep}_2((C_3)^4)=7},
\qquad
\boxed{\operatorname{sep}_2((C_5)^3)=7}.
\]

\subsubsection{Computational proof}
The lower bounds are respectively seven, since \(2^6<81\le2^7\), and seven,
since \(2^6<125\le2^7\). The repository stores explicit Boolean colourings and
seven-element windows in \texttt{c3four\_size7\_certificate.json} and
\texttt{c5cubed\_size7\_certificate.json}. The generic direct verifier checks
that the translated observation words are pairwise distinct: 81 words in the
first case and 125 in the second. Thus both upper bounds equal the counting
lower bounds. \(\square\)

\subsection{6. Groups of order 32}
There are 51 groups of order \(32\) in the GAP SmallGroups library.

\subsubsection{Proposition 6.1}
For every group \(G\) of order \(32\),
\[
\boxed{\operatorname{sep}_2(G)=5.}
\]
\subsubsection{Proof}
The counting bound gives \(\operatorname{sep}_2(G)\ge5\). For each group \(\operatorname{SmallGroup}(32,i)\), \(1\le i\le51\), the repository stores an explicit five-element window and a Boolean colouring. The verification script reconstructs the multiplication table of each group in GAP and checks directly that the 32 translated words are pairwise distinct. Hence \(\operatorname{sep}_2(G)\le5\). Equality follows. \(\square\)

\subsection{7. Reproducibility}
The repository contains the finite verifiers, SAT search programs, orbit reduction for \((C_3)^3\), sharp certificates for \((C_3)^4\) and \((C_5)^3\), and explicit certificates for all 51 groups of order \(32\). The generic search tool treats an exhausted bounded search as inconclusive rather than as a proof of nonexistence.

Ming-Hsuan Kang independently reproduced the original core computational package.

\subsection{8. Open problems}
\begin{enumerate}
\item Classify finite groups \(G\) satisfying \(\operatorname{sep}_2(G)=\lceil\log_2|G|\rceil\).
\item Determine how large \(\delta_q(G)=\operatorname{sep}_q(G)-\lceil\log_q|G|\rceil\) can be.
\item Find structural conditions forcing \(\delta_q(G)=0\).
\end{enumerate}
The example \((C_3)^3\) shows that universal equality with the information lower bound is false, while the order-32 computation shows that noncommutativity alone does not force positive defect.

\subsection{Acknowledgments}
I thank Ming-Hsuan Kang and Yu-Hsuan Hsieh for helpful correspondence concerning their separating-pattern parameter. I also thank Ming-Hsuan Kang for independently reproducing the original core computational package.

\subsection{References}
\begin{enumerate}
\item M.-H. Kang and Y.-H. Hsieh, \emph{Information and Locality in Cayley Graphs}, arXiv:2608.04608v1, 2026.
\item B. Cameron, A. Gündoğan, and J. Sawada, \emph{Cut-Down de Bruijn Sequences}, arXiv:2205.02815, 2022.
\end{enumerate}

\end{document}
